\chapter{Топологические и метрические пространства}

\lesson{1}{от 4 сен 2024 12:45}{Начало}


\section{Топологические пространства}

\subsection{Основные понятия}

Пусть $X$ -- непустое множество, через $\mathcal{P}(X)$ будем обозначать совокупность всех его подмножеств (часто для этой цели используется обозначение $2^X$). Подмножества $\mathcal{P}(X)$ в зависимости от стилистического удобства иногда будем называть совокупностями или классами.

\begin{definition}[Топологическое пространство (ТП)]
	Множество $X$ называется \emph{топологическим пространством (ТП)}, если на нем определено отображение
	\[
		O: \ X \ni x \rightarrow O(x) \subset \mathcal{P}(X)
	\]
	такое, что для любого $x \in X$ класс $O(x)$ обладает следующими свойствами:
	\begin{enumerate}
		\item $\forall V \in O(x) \ x \in V$.
		\item $\forall U,V \in O(x) \ U \cap V \in O(x)$.
		\item $\forall  U \in O(x) \ \big(U \subset V \subset X \implies V \in O(x)\big)$.
		\item $\forall V \in O(x) \ \exists U \in O(x): \ U \subset V \ \land \big(\forall y \in U \ U \in O(y)\big)$.
	\end{enumerate}

	Элементы класса $O(x)$ называются окрестностями точки $x \in X$.
\end{definition}

\begin{remark}
	В дальнейшем множество всех окрестностей точки $x \in X$ в топологическом пространстве $X$ вместо $O(x)$ для кратности будем обозначать $O_{x}^{X} $ или просто $O_x$.
\end{remark}

\begin{definition}[Окрестность множества]
	Подмножество топологического пространства $X$ называется \emph{окрестностью множества} $A \subset X$, если оно является окрестностью каждой точки множества $A$.
\end{definition}

\begin{definition}[Проколотая окрестность]
	\emph{Проколотой окрестностью} точки топологического пространства называется множество, полученной удалением самой этой точки из некоторой ее окрестности, например, если $V$ -- окрестность точки $x$ в топологическом пространстве $X$, то $\dot{V} := V \setminus \{x\} $ -- проколотая окрестность точки $x$.
\end{definition}

\begin{definition}[Открытое подмножество]
	Подмножество топологического пространства называется \emph{открытым}, если оно является окрестностью каждой своей точки. Совокупность всех открытых подмножеств топологического пространства $X$ обозначается
	\[
		Op \ X.
	\]
\end{definition}

\begin{remark}
	Согласно определению топологического пространства окрестности его точек обладают следующими свойствами:
	\begin{enumerate}
		\item Каждая окрестность точки содержит эту точку.
		\item Пересечение конечного числа окрестностей точки также является окрестностью этой точки.
		\item Каждое множество, содержащее окрестность точки, является окрестностью этой точки.
		\item В каждой окрестности точки содержится открытая окрестность этой точки.
	\end{enumerate}

	Перефразируя определение топологического пространства, можно сказать, что топологическое пространство -- это непустое множество, каждому элементу ....
\end{remark}

\begin{definition}[Замкнутое подмножество]
	Подмножество топологического пространства называется \emph{замкнутым}, если его дополнение открыто. Совокупность всех замкнутых подмножеств топологического пространства $X$ обозначается
	\[
		Cl \ X.
	\]
\end{definition}

\begin{remark}
	Нетрудно доказать, что замкнутые подмножества топологического пространства обладают следующими свойствами:
	\begin{enumerate}
		\item Пустое множество и все пространство замкнуты.
		\item Объединение конечного числа замкнутых множеств замкнуто.
		\item Пересечение любой совокупности замкнутых множеств замкнуто.
	\end{enumerate}
\end{remark}

\begin{remark}
	Имея топологическое пространство, мы имеем множество, на котором задана топологическая структура (для кратности удобнее говорить \emph{топология}), то есть множества окрестностей всех его точек, все его открытые и все замкнутые подмножества. Причем, как нетрудно заметить, для задания топологической структуры на множестве не важно с чего начинать, с окрестностей, с открытых или замкнутых множеств. ....
\end{remark}

\begin{definition}[Подпространство ТП]
	Если $X$ -- топологическое пространство, а $X_0$ -- его непустое подмножество, то $X_0$ мы также можем рассматривать как топологическое пространство, наделив его \emph{индуцированной} из $X$ топологией. Эта топология определеяется следующим образом:
	\[
		\forall x \in X_0 \ O_{x}^{X_0} \coloneq \big\{V \cap X_0 \ \big| \ V \in O_{x}^{X} \big\},
	\]
	\[
		Op \ X_0 \coloneq \big\{A \cap X_0 \ \big| \ A \in Op \ X\big\}, \qquad Cl \ X_0 \coloneq \big\{B \cap X_0 \ \big| \ B \in Cl \ X\big\}.
	\]

	При этом $X_0$ называется \emph{подпространством топологического пространства} $X$. Согласно сказанному в предыдущем замечании для определения индуцированной топологии достаточно задания одного из трех приведенных классов множеств.
\end{definition}

\begin{definition}
	Если на множестве заданы две топлогические структуры (топлогия), то говорят, что вторая \emph{мажорирует} первую, если каждое множество, открытое в первой топологии, будет открыто и во второй, то есть множесв, открытых во второй топлогии, теоретико-множественно не меньше, чем открытых в первой.

	Это равносильно такому же соотношению окрестностей точек в этих топологиях. При этом также говорят, что вторая топологическая структура (топология) \emph{не слабее} первой, а первая \emph{не сильнее} второй.
\end{definition}

\begin{remark}
	Когда на одном и том же множестве рассматривают несколько топологических структур, то для обозначения соответсвующих совокупностей открытых множеств (их обычно называют топлогиями) используют буквенные обозначения типа $\sigma , \ \tau , \ v$ и т.д. При этом множество $X$, наделенное топологией $\tau $, обозначается
	\[
		(X,\tau ).
	\]
\end{remark}

\begin{definition}[База окрестностей]
	Совокупность $V \subset O_{x}^{X} $ называется \emph{базой окрестностей} точки $x \in X$, если в каждой окрестности точки $x$ содержится некоторое множество из $V$. Например, в обычной топологии на числовой прямой каждая точка $x$ имеет базу окрестностей, состоящую из так называемых $\epsilon $-окрестностей точки $x$, то есть интервалов вида $(x - \epsilon ; \ x + \epsilon ) \ (\epsilon > 0)$.
\end{definition}

\begin{definition}[ТП, удовл. первой аксиоме счетности]
	Говорят, что топологическое пространство \emph{удовлетворяет первой аксиоме счетности}, если каждая его точка имеет не более чем счетную базу окретсностей.
\end{definition}

\begin{definition}[ТП, удовлю второй аксиоме счетности]
	Говорят, что топологическое пространство \emph{удовлетворяет второй аксиоме счетности}, если существует не более чем счетная совокупность его открытых подмножеств такая, что любое его непустое открытое множество является объединением каких-то элементов этой совокупности.
\end{definition}

\begin{remark}
	Нетрудно заметить, что топологическое пространство, удовлетворяющее второй аксиоме счетности, удовлетворяет и первой аксиоме счетности. Простейшим примером пространства, удовлетворяющего второй аксиоме, является ....
\end{remark}

\begin{definition}[Внутренняя точка]
	Точка $x \in A$ называется \emph{внутренней точкой множества} $A$, если $A$ является окрестностью этой точки. Множество всех внутренних точек множества называется внутренностью и обозначается
	\[
		int \ A.
	\]
\end{definition}

\begin{definition}[Точка прикосновения]
	Точка $x \in X$ называется \emph{точкой прикосновения} множества $A$, если пересечение любой ее окрестности с множеством $A$ не пусто. Множество всех точек прикосновения множества $A$ называется \emph{замыканием} множества и обозначается
	\[
		cl \ A.
	\]
\end{definition}

\begin{definition}[Предельная точка]
	Точка называется \emph{предельной точкой} подмножества топологического пространства, если каждая окрестность этой точки содержит точки множества, отличные от данной.
\end{definition}

\begin{definition}[Граничная точка]
	Точка $x \in X$ называется \emph{граничной точкой множества} $A$, если любая ее окрестность имеет непустое пересечение как с множеством $A$, так и с его дополнением. Множество всех граничных точек множества $A$ называется его границей и обозначается
	\[
		\partial A.
	\]
\end{definition}
